\section{Résultats}

\subsection{Position des hacktivistes face à la désinformation}

L'étude révèle un consensus parmi les hacktivistes : aucun d'entre eux n'accueille favorablement l'appropriation néfaste du trolling et des "mêmes" pour la propagande politique. Fidèles à l'éthique hacker originelle (A.K.A. RMS), ils perçoivent la désinformation comme une menace pour la vision démocratique d'Internet.

\begin{itemize}
	\item \textbf{Conflit avec l'éthique hacker} : La désinformation contredit le postulat selon lequel « toute information doit être libre » en créant un désordre informationnel qui catalyse la polarisation.
	\item \textbf{Méfiance envers l'autorité} : Elle s'oppose au principe de décentralisation car elle est souvent orchestrée par une « autorité fantôme » ou étatique.
	\item \textbf{Détournement des outils} : Les participants notent que les outils qu'ils ont popularisés ("mêmes", trolling) sont désormais retournés contre l'éthique hacker elle-même pour manipuler les émotions.
	\item \textbf{Critique des plateformes} : Les hacktivistes accusent les réseaux sociaux d'être les facilitateurs de ce désordre en monétisant l'engagement plutôt qu'en appliquant des politiques strictes de participation.
\end{itemize}

\subsection{Méthodes préconisées}


Pour contrer ce développement, les hacktivistes recommandent d'utiliser des méthodes éprouvées pour protéger le discours social.

\begin{itemize}
	\item \textbf{Actions de neutralisation} : Les interventions privilégiées incluent le \textit{deplatforming} (fermeture de comptes), le \textit{doxing} (révélation d'identité) des financeurs et le \textit{leaking} de documents.
	\item \textbf{Saturation informationnelle} : Ils préconisent de combattre la désinformation par un surplus d'informations véridiques, une stratégie de « combattre le feu par le feu ».
	\item \textbf{L'exemple de \#OpJane} : Cette opération illustre l'usage de « désinformation plausible » pour saturer et rendre inutilisables les bases de données de « chasseurs de primes » anti-avortement au Texas.
	\item \textbf{Pression économique} : Le \textit{doxing} est également suggéré pour inciter les annonceurs à retirer leur soutien aux influenceurs propageant des faussetés.
\end{itemize}

\subsection{Recommandations pour la littératie informationnelle}

Au-delà des actions techniques, la majorité des participants recommande des interventions éducatives.

\begin{itemize}
	\item \textbf{Réforme éducative} : Ils appellent à l'intégration de programmes de pensée critique dans les écoles pour pallier les lacunes actuelles.
	\item \textbf{Responsabilité numérique} : L'enseignement des compétences de sécurité opérationnelle (OpSec) devrait être vu comme une responsabilité sociale.
	\item \textbf{Bots de vérité} : Ils proposent le développement de « bots de vérité » pour faire face aux bots de désinformation et aider les utilisateurs à localiser les faits.
	\item \textbf{Neutralité du ton} : Ils estiment que la littératie actuelle échoue car elle adopte un ton de « cancel culture » plutôt qu'une attitude impartiale rappelant que les faits scientifiques n'ont pas de propriétés politiques.
\end{itemize}
